\begin{resumo}[Abstract]
\begin{otherlanguage*}{english}
Large Language Models (LLMs) have been increasingly used in software development, including tasks such as code generation, explanation, and correction. However, their use in software security still requires caution, since automatically generated fixes may be incomplete, insecure, or may compromise the original behavior of the application. This work investigates the integration of LLMs, Static Application Security Testing (SAST) tools, automated tests, and contextualization strategies to support the correction of security issues in source code. The proposed methodology applies an automated pipeline to an open source project, using the vulnerabilities effectively detected by the SAST tool. As an initial step, a proof of concept was developed to validate the workflow for issue identification, fix generation, static reanalysis, test execution, and metric consolidation. Thus, this work establishes the technical and methodological foundations to evaluate the effectiveness of the pipeline in TCC 2.

\textbf{Keywords}: large language models; static application security testing; vulnerability repair; prompt engineering; RAG; pipeline.
\end{otherlanguage*}
\end{resumo}
